% D06-001. PDR-05, §6.8, §11 и обсуждение.
\section{Инженерная проверка результата модернизации}
\label{sec:ch5_engineering_verification}
\conceptsource{Статья 5, плотная оценка и программа RTDS.
Перенесена логика проверки полного состояния, ошибок
во времени и переходов. Окончательный рейтинг
незавершённой статьи здесь не утверждается.}

(Примечание автора: Ну наверное стоит писать не «здесь проверяется», а в этом разделе, или чёт такое, всё же это диссертация. Ну или «Далее описана проверка того-то».)

Здесь проверяется поведение органа на записи во времени.
Важно установить не только долю правильных ответов,
но и длительность ошибки, лишние переключения и реакцию
на изменение режима. Например, одинаковый процент совпадений
может соответствовать коротким помеховым ответам или длительной
неверной блокировке. Для работы БАВР эти ситуации различаются.


\subsection{Полное состояние и последствия разных ошибок}
\conceptsource{Статья 5, подраздел 6.8. Три состояния,
раздельные знаменатели и ошибочное время.}

\iffalse
(Примечание автора: Изменения характеристики? Какой характеристики? Не понял.)
(Примечание автора: Ок, закрываем задачу.)
\fi
\authorreviewnote{AA167}

После изменения характеристики проверяется
не только средняя доля совпавших направлений.
Полный ответ имеет три состояния.
Направление может быть прямым, обратным
либо не определяться в данном контексте.
При неприменимости знак направления не используется
как содержательное решение.

Основной показатель описывает суммарное время
несовпадения полного состояния с экспертным ответом.
Оно относится к длительности наблюдения
и может выражаться в миллисекундах на секунду.
Например, 30~мс/с означают 3\% ошибочного времени.
Такое представление удобно для осциллограмм
с разной длительностью и частотой регистрации.

Для определения причины общий показатель дополняется
условными долями. Ложное прямое направление
рассчитывается на времени экспертного обратного,
ложное обратное на времени прямого.
Лишний отказ относится ко времени применимости,
а пропуск неприменимости к неприменимому времени.
Эти величины имеют разные знаменатели
и не складываются между собой.
Если нужного состояния в записи нет,
соответствующая доля не определена.

\iffalse
(Примечание автора: Чёт выше как-то заумно написано, хорошо бы немного упростить.)
(Примечание автора: Ок, закрываем задачу.)
\fi
\authorreviewnote{AA168}

В логике БАВР ошибочное прямое состояние
может сохранять блокирующий признак при выбеге.
Ошибочное обратное может снимать его
при повреждении на стороне нагрузки.
Но число таких отсчётов ещё не равно числу
ошибочных переключений оборудования.
Окончательное действие определяется
совокупностью внешних условий автоматики.

\subsection{Длительные ошибки и устойчивость выхода}
\conceptsource{Статья 5, непрерывная ошибка, неблагоприятное
окно и плотность лишних переключений.}

Одинаковая средняя доля может получаться
при множестве коротких несовпадений
или при одной продолжительной ошибке.
Поэтому определяется максимальная непрерывная ошибка
и число записей с участками длиннее
5, 20, 50 и 100~мс.
Дополнительно находится ошибочное время
в наиболее неблагоприятном окне 100~мс.
Последний показатель сохраняет чувствительность
к продолжительной ошибке, даже если она
прерывается короткими правильными ответами.
Указанные интервалы служат диагностическими
ориентирами исследования, а не нормативами БАВР.

\iffalse
(Примечание автора: К длительному неблагополучию… Пиши более внятно, пожалуйста, буду повторять, не надо так заковыристо. Да и разделы местами как будто мелкие слишком, зачем их в подразделы-то выделять, как этот например.)
(Примечание автора: Ок, закрываем задачу.)
\fi
\authorreviewnote{AA169}

Частота лишних переключений рассматривается
вне начальной неготовности и экспертных переходных зон.
Учитываются изменения направления
и применимости. Уменьшение дребезга само по себе
не означает улучшения, если орган стал дольше
удерживать неправильное состояние.
Поэтому устойчивость рассматривается вместе
с длительностью и физическим содержанием ошибки.

\subsection{Реакция на смену экспертного состояния}
\conceptsource{Статья 5, сопоставление однотипных переходов
и удержание нового состояния. Пропуски остаются в оценке.}

Переходы между направлениями сопоставляются
отдельно от потери и восстановления применимости.
Одно изменение выхода органа может соответствовать
только одному экспертному переходу того же типа.
В исследовании используется окно сопоставления
$\pm100$~мс и дополнительная проверка
чувствительности при $\pm50$~мс.
Сначала выбирается наибольшее число сопоставимых пар,
затем наименьший суммарный временной сдвиг.

Положительный сдвиг означает запаздывание,
отрицательный опережение экспертной границы.
Новое состояние дополнительно проверяется
на удержание в течение 0, 5 и 20~мс.
Это условие анализа полученного ответа,
а не добавленная в измерительный орган выдержка.
При опережении состояние должно сохраняться
до экспертной границы и после неё.

Медиана только найденных переходов
не описывает пропущенные события.
Рядом с ней определяется доля всех достаточно
наблюдаемых экспертных переходов, на которые
получен своевременный и удержанный ответ.
Завершение записи и слишком короткое новое
экспертное состояние учитываются отдельно.
Они отличаются от пропуска при полностью
наблюдаемом временном интервале.

\subsection{Проверка на связанных режимах БАВР}
\conceptsource{Статья 5, раздел 11. Сто RTDS-опытов,
две секции и отдельная проверка редкой неприменимости.}

\iffalse
(Примечание автора: можно бы добавить, что это выполнено по программе тестирования Транснефти, или смежной с ней… Как-то обыграть это (или вообще не называть Транснефть, прост программа испытаний). Ибо в реальности, это и есть прям программа, которую я потом доразметил вручную, но испытания прям реально-требуемые, как минимум сами процессы переключений.)
(Примечание автора: Ок, закрываем задачу.)
\fi
\authorreviewnote{AA170}

Архивная статистика дополняется предметными
испытаниями, связанными с назначением БАВР.
В исходной работе использованы 100 RTDS-опытов
с двумя контролируемыми секциями.
Каждая запись длится 3,2~с
при частоте регистрации 20~кГц.
Две секции одного опыта электрически связаны,
поэтому 12,8 млн отсчётов их состояния
не являются таким же числом независимых испытаний.

Показатели сначала определяются по секциям,
затем по опытам. Для максимальной ошибки
опыта используется худшая секция.
Источники реальных записей рассматриваются раздельно.
Так средняя оценка по удобному массиву
не скрывает ухудшения на целевых режимах.

Для направления вместе с совпадением
указывается покрытие экспертно применимого времени.
Это показывает, не достигнуто ли высокое качество
на малой части удобных отсчётов.
Отдельно проверяется обнаружение неприменимости.
Редкость такого состояния может давать
высокий общий процент совпадений даже органу,
который его ни разу не выделил.

\iffalse
(Примечание автора: Ещё раз повторю, ты как-то на подглавы разбиваешь всё подряд, может где-то и нужно, но не по пару абзацев же? Ну как-то уж очень мелок… Где-то может текста не дописала, а где-то, может и объединить можно было.)
(Примечание автора: Ок, закрываем задачу.)
\fi
\authorreviewnote{AA171}

\subsection{Обратная связь с изменением органа}
\conceptsource{Статья 5, обсуждение адаптации под объект;
обобщение автора 27.10.2026 о формульной и нейросетевой
реализации одной методики.}

Завершающий результат проверки имеет
предметное содержание. Устанавливается,
в каких режимах прежняя характеристика
недостаточна и какое изменение устраняет
наблюдаемую причину.
В аналитическом органе могут изменяться
опорная величина, память, рабочая область
и зависимость порога от режима.
В нейросетевом варианте корректируются
данные, цели и обучаемое преобразование.

После локальной коррекции повторяется
проверка на сложных и обычных режимах.
Устранение одного неправильного ответа
не должно оцениваться без рассмотрения
остальных условий работы.
Для настройки под конкретный объект
сохраняется отдельный проверочный материал,
не использованный при этой настройке.

Таким образом, на примере РНМ рассматривается
общий способ совершенствования органа РЗА.
Он связывает физическое назначение,
измерительный контекст, массовую статистику,
экспертный разбор и повторную инженерную проверку.
Выбор математической реализации входит
в этот процесс, а не определяет его цель заранее.

